% docs/slides/preamble.tex -- 五份 deck 共享
\documentclass[aspectratio=169,11pt]{beamer}

% ==== 中文与字体 ====
\usepackage[UTF8,fontset=none]{ctex}
\usepackage{fontspec}

\IfFontExistsTF{PingFang SC}{%
  \setCJKmainfont{PingFang SC}[BoldFont=PingFang SC,ItalicFont=PingFang SC,AutoFakeSlant=0.15]
  \setCJKsansfont{PingFang SC}[ItalicFont=PingFang SC,AutoFakeSlant=0.15]
  \setCJKmonofont{PingFang SC}
}{%
  \IfFontExistsTF{Source Han Sans SC}{%
    \setCJKmainfont{Source Han Sans SC}
    \setCJKsansfont{Source Han Sans SC}
  }{%
    \IfFontExistsTF{Noto Sans CJK SC}{%
      \setCJKmainfont{Noto Sans CJK SC}
      \setCJKsansfont{Noto Sans CJK SC}
    }{%
      \setCJKmainfont{Heiti SC}
      \setCJKsansfont{Heiti SC}
    }
  }
}

\IfFontExistsTF{Helvetica Neue}{\setmainfont{Helvetica Neue}}{%
  \IfFontExistsTF{Inter}{\setmainfont{Inter}}{}}
\IfFontExistsTF{JetBrains Mono}{\setmonofont{JetBrains Mono}[Scale=0.85]}%
  {\setmonofont{Menlo}[Scale=0.85]}

% ==== 颜色 ====
\usepackage{xcolor}
\definecolor{cMain}{HTML}{1F4E79}
\definecolor{cAccent}{HTML}{E07A1F}
\definecolor{cGray}{HTML}{595959}
\definecolor{cMute}{HTML}{F4F4F4}
\definecolor{cOK}{HTML}{2E7D32}
\definecolor{cBad}{HTML}{B23A3A}
\definecolor{cWarn}{HTML}{C8A300}

% ==== 主题 ====
\usepackage{tcolorbox}
\usetheme{metropolis}
\metroset{progressbar=frametitle,numbering=fraction,block=fill}
\setbeamercolor{frametitle}{bg=cMain,fg=white}
\setbeamercolor{progress bar}{fg=cAccent,bg=cMute}
\setbeamercolor{alerted text}{fg=cAccent}
\setbeamercolor{title}{fg=cMain}

% metropolis 默认尝试 Fira Sans，未安装时 fallback 到 Latin Modern Sans。
% 这里在主题加载后显式覆盖，强制使用 macOS 上一定存在的 Helvetica Neue，
% 跟中文 PingFang 风格一致；找不到则按顺序尝试 Inter / Arial。
\IfFontExistsTF{Helvetica Neue}{%
  \setsansfont{Helvetica Neue}%
}{%
  \IfFontExistsTF{Inter}{\setsansfont{Inter}}{%
    \IfFontExistsTF{Arial}{\setsansfont{Arial}}{}}%
}

% ==== TikZ ====
\usepackage{tikz}
\usetikzlibrary{positioning,arrows.meta,fit,shapes.geometric,calc,%
  decorations.pathmorphing,backgrounds}

\tikzset{
  node-box/.style={draw=cMain,thick,rounded corners=2pt,
    minimum height=8mm,minimum width=18mm,align=center,font=\small},
  node-state/.style={draw=cMain,thick,circle,minimum size=14mm,
    align=center,font=\small},
  node-ok/.style={draw=cOK,thick,fill=cOK!10},
  node-warn/.style={draw=cWarn,thick,fill=cWarn!15},
  node-bad/.style={draw=cBad,thick,fill=cBad!10},
  % 色盲友好的状态机三态：蓝 (正常/Closed) / 橙 (告警/Open) / 灰 (中间/HalfOpen)
  % 避免 red-green 二元区分，改用 hue + lightness 双通道。
  node-state-normal/.style={draw=cMain,thick,fill=cMain!12},
  node-state-alert/.style={draw=cAccent,thick,fill=cAccent!18},
  node-state-transition/.style={draw=cGray,thick,fill=cGray!18},
  arrow/.style={-{Stealth[length=2mm]},thick,draw=cMain},
  arrow-async/.style={-{Stealth[length=2mm]},thick,dashed,draw=cGray},
  arrow-bad/.style={-{Stealth[length=2mm]},thick,draw=cBad},
  arrow-alert/.style={-{Stealth[length=2mm]},thick,draw=cAccent},
  lifeline/.style={dashed,thick,draw=cGray},
}

% ==== 代码 ====
\usepackage{listings}
\lstdefinelanguage{Go}{
  morekeywords={break,case,chan,const,continue,default,defer,else,
    fallthrough,for,func,go,goto,if,import,interface,map,package,range,
    return,select,struct,switch,type,var,nil,true,false,iota},
  morekeywords=[2]{string,int,int64,uint,uint64,bool,byte,error,context},
  sensitive=true,
  morecomment=[l]//,morecomment=[s]{/*}{*/},
  morestring=[b]",morestring=[b]`,
}
\lstdefinelanguage{LuaRedis}{
  morekeywords={and,break,do,else,elseif,end,false,for,function,if,in,
    local,nil,not,or,repeat,return,then,true,until,while},
  morekeywords=[2]{redis,call,KEYS,ARGV,tonumber,HGET,HSET,GET,SET,
    DECR,DECRBY,INCR,INCRBY,EXPIRE,PEXPIRE,ZADD,ZCARD,ZREMRANGEBYSCORE,
    EVAL,EVALSHA},
  sensitive=true,
  morecomment=[l]{--},
  morestring=[b]",morestring=[b]',
}

\lstdefinestyle{base}{
  basicstyle=\ttfamily\scriptsize,
  numbers=left,numberstyle=\tiny\color{cGray},
  keywordstyle=\color{cMain}\bfseries,
  keywordstyle=[2]\color{cAccent},
  stringstyle=\color{cOK},
  commentstyle=\color{cGray}\itshape,
  backgroundcolor=\color{cMute},
  frame=single,framerule=0pt,
  showstringspaces=false,
  breaklines=true,
  tabsize=2,
  columns=fullflexible,
}
\lstset{style=base}

% ==== 三线表与附录页码 ====
\usepackage{booktabs}
\usepackage{appendixnumberbeamer}

% ==== 页脚来源标注 ====
\newcommand{\srcnote}[1]{{\color{cGray}\scriptsize\itshape #1}}

% 关键论点框
\newcommand{\keypoint}[1]{%
  \begin{tcolorbox}[colback=cMute,colframe=cMain,boxrule=0.5pt,arc=2pt]%
  #1\end{tcolorbox}}


\title{商品详情页业务实践}
\subtitle{从 \texttt{product} 表字段到 Cache Aside 的业务取舍}
\author{gomall · 业务侧扩展}
\date{}

\begin{document}

\maketitle

\begin{frame}{今日大纲}
\tableofcontents
\end{frame}

% ============== 1. 业务地位 ==============
\section{商品详情页的业务地位}

\begin{frame}{从流量入口看详情页}
\begin{tikzpicture}[node distance=5mm and 14mm]
  \node[node-box] (home) {首页推荐};
  \node[node-box,below=of home] (list) {分类列表};
  \node[node-box,below=of list] (search) {搜索结果};
  \node[node-box,below=of search] (rec) {个性化推荐};
  \node[node-box,below=of rec] (share) {社交分享};
  \node[node-box,right=22mm of search,node-warn,minimum width=26mm] (detail) {商品详情页\\\scriptsize\texttt{/product/show}};
  \node[node-box,right=18mm of detail,node-ok] (cart) {加购 / 下单};
  \draw[arrow] (home.east) -- (detail.west);
  \draw[arrow] (list.east) -- (detail.west);
  \draw[arrow] (search.east) -- (detail.west);
  \draw[arrow] (rec.east) -- (detail.west);
  \draw[arrow] (share.east) -- (detail.west);
  \draw[arrow] (detail) -- node[above,font=\tiny]{转化} (cart);
\end{tikzpicture}
\vspace{2mm}
\begin{itemize}
  \item 五条入口路径汇聚到\textbf{同一个接口} —— 任何一条爆量都会引发读放大。
  \item 详情页是转化漏斗里离"下单"最近的一步：延迟 / 不一致直接砸 GMV。
  \item 业界长期观察：详情页首屏每多一秒，跳出率显著上升、加购率显著下降。
\end{itemize}
\srcnote{service/product.go:38--68 \texttt{ProductShow}}
\end{frame}

% ============== 2. product 表字段 ==============
\section{\texttt{product} 表字段的业务含义}

\begin{frame}[fragile]{gomall 的 \texttt{Product} 模型}
\begin{lstlisting}[language=Go,firstnumber=12,
  caption={\srcnote{repository/db/model/product.go:12--26}}]
type Product struct {
    gorm.Model
    Name          string `gorm:"size:255;index"`  // 商品短名
    CategoryID    uint   `gorm:"not null"`         // 一级分类
    Title         string                            // 长标题/卖点
    Info          string `gorm:"size:1000"`         // 详情正文
    ImgPath       string                            // 主图
    Price         string                            // 原价
    DiscountPrice string                            // 折后价
    OnSale        bool   `gorm:"default:false"`     // 在售开关
    Num           int                                // 库存数（DB 真值）
    BossID        uint                               // 商家 ID
    BossName      string                             // 冗余商家名
    BossAvatar    string                             // 冗余商家头像
}
\end{lstlisting}
\end{frame}

\begin{frame}{Name / Title / Info / Price / DiscountPrice}
\begin{tikzpicture}[node distance=4mm and 10mm]
  \node[node-box,minimum width=22mm] (n) {\textbf{Name}\\\scriptsize 列表卡片};
  \node[node-box,minimum width=22mm,right=of n,node-warn] (t) {\textbf{Title}\\\scriptsize 首屏 H1};
  \node[node-box,minimum width=22mm,right=of t,node-ok] (i) {\textbf{Info}\\\scriptsize 正文 1000 字};
  \draw[arrow] (n) -- node[above,font=\tiny]{信息密度 $\uparrow$} (t);
  \draw[arrow] (t) -- (i);
\end{tikzpicture}
\vspace{2mm}
\begin{itemize}
  \item \texttt{Name} 上加 \texttt{index}：列表 / 搜索靠它；\texttt{Title} 承担首屏注意力；\texttt{Info} 1000 字决定停留时长但\textbf{不进索引}。
  \item \texttt{Price} 是\textbf{划线价}，\texttt{DiscountPrice} 是真实结算价 —— 制造\textbf{价格锚点} + 支持大促只动后者。
  \item 两个金额字段都是 \texttt{string}，\textbf{不能}直接算术：金额一律走 \texttt{decimal} / 整数分。
\end{itemize}
\srcnote{repository/db/model/product.go:14--20}
\end{frame}

\begin{frame}{Num — DB 真值与 Redis 预占的关系}
\begin{tikzpicture}[node distance=6mm and 14mm]
  \node[node-box,node-ok] (db) {\texttt{product.Num}\\\scriptsize DB 真值};
  \node[node-box,right=of db,node-warn] (redis) {\texttt{stock:\{pid\}}\\\scriptsize Redis 可售余量};
  \node[node-box,right=of redis] (user) {下单请求};
  \draw[arrow] (user) -- node[above,font=\tiny]{Lua DECR} (redis);
  \draw[arrow-async] (redis) -- node[above,font=\tiny]{异步落库} (db);
\end{tikzpicture}
\vspace{2mm}
\begin{itemize}
  \item 详情页展示的"库存剩 N"是\textbf{展示值}，可走 DB 也可走缓存。
  \item 真正"能不能下单"在 Redis 预占（见 Deck 02），\texttt{Num} 是\textbf{对账依据}：Saga 失败、订单超时取消都回滚到这里。
  \item \textbf{结论}：详情页缓存里写 \texttt{Num} 没问题，下单链路必须读\textbf{运行时余量}。
\end{itemize}
\srcnote{repository/db/model/product.go:22}
\end{frame}

\begin{frame}{OnSale / CategoryID / Boss* 三组字段}
\begin{itemize}
  \item \texttt{OnSale}：商家自助、平台风控、运营批量、算法\textbf{都能改}；普通用户、缓存层\textbf{不能改}。
  \item \texttt{CategoryID}：承担 \textbf{SEO}（关键词 / 面包屑）+ \textbf{推荐}（"你可能感兴趣"）+ \textbf{业务分组}（限购 / 退换 / 税率）三重职责，所以加了 \texttt{not null}。
  \item \texttt{BossID / BossName / BossAvatar}：冗余\textbf{不是}设计错误：
  \begin{itemize}
    \item 详情页一次查询展示商家头像 + 名称 —— 避免每次 JOIN 商家表。
    \item 商家改名 / 换头像\textbf{低频}，可接受短暂不同步。
    \item 整页缓存友好：1 次 GET 给前端全部素材。
  \end{itemize}
\end{itemize}
\keypoint{冗余是"读多写少"业务的工程妥协，不是范式违规。}
\srcnote{repository/db/model/product.go:15, 21, 23--25}
\end{frame}

% ============== 3. 谁能更新 ==============
\section{谁能更新商品}

\begin{frame}{四类写入主体共享同一条写路径}
\begin{tikzpicture}[node distance=5mm and 14mm]
  \node[node-box] (merchant) {商家自助};
  \node[node-box,below=of merchant] (risk) {平台风控};
  \node[node-box,below=of risk] (ops) {运营批量};
  \node[node-box,below=of ops] (algo) {算法降价};
  \node[node-box,right=18mm of risk,node-warn,minimum width=28mm] (api) {\texttt{ProductUpdate}\\\scriptsize service/product.go:271};
  \node[node-box,right=16mm of api,node-ok] (db) {DB + 双删};
  \draw[arrow] (merchant.east) -- (api);
  \draw[arrow] (risk.east) -- (api);
  \draw[arrow] (ops.east) -- (api);
  \draw[arrow] (algo.east) -- (api);
  \draw[arrow] (api) -- (db);
\end{tikzpicture}
\vspace{1mm}

\begin{tabular}{@{}llll@{}}
\toprule
主体 & 频率 & 改的字段 & 一致性诉求 \\
\midrule
商家自助 & 中 & \texttt{Title/Info/Price/...} & 自己看得到 \\
平台风控 & 低 & \texttt{OnSale=false} & 立即下架 \\
运营批量 & 突发 & \texttt{DiscountPrice/OnSale} & 大促瞬间生效 \\
算法降价 & 高 & \texttt{DiscountPrice} & 分钟级容忍 \\
\bottomrule
\end{tabular}
\srcnote{service/product.go:271--292}
\end{frame}

% ============== 4. 延迟的业务影响 ==============
\section{更新延迟的业务影响}

\begin{frame}{事故 1：下架后用户仍能买}
\begin{tikzpicture}[node distance=8mm]
  \foreach \name/\x in {商家/0, DB/30, Redis/60, 用户 A/95, 用户 B/130} {
    \node[node-box,minimum width=14mm] at (\x mm,0) (\name) {\name};
    \draw[lifeline] (\x mm,-2mm) -- (\x mm,-44mm);
  }
  \draw[arrow] (0,-8mm) -- node[above,font=\tiny]{下架} (30mm,-8mm);
  \draw[arrow-bad] (95mm,-16mm) -- node[above,font=\tiny\color{cBad}]{读旧缓存} (60mm,-16mm);
  \draw[arrow-bad] (130mm,-24mm) -- node[above,font=\tiny\color{cBad}]{读旧缓存} (60mm,-24mm);
  \draw[arrow-bad] (95mm,-32mm) -- node[above,font=\tiny\color{cBad}]{下单成功} (30mm,-32mm);
  \draw[arrow-bad] (130mm,-40mm) -- node[above,font=\tiny\color{cBad}]{付款} (30mm,-40mm);
\end{tikzpicture}
\vspace{1mm}
\small 商家已下架但缓存未失效 —— 用户付完钱才发现卖家不发货，客服压力暴涨。
\end{frame}

\begin{frame}{事故 2：改价后旧价单的业务处理}
\begin{itemize}
  \item 大促 0:00:00，运营把 \texttt{DiscountPrice} 从 ¥199 改到 ¥99。
  \item 0:00:00.300，缓存还没失效，用户 A 看到 ¥199 下单。业务部门要选：
  \begin{itemize}
    \item \textbf{让利方案}：按 ¥99 结算，多付差价退还（最常见）。
    \item \textbf{补差方案}：保持 ¥199，下发"漏发优惠券"安抚。
    \item \textbf{严格方案}：让 A 重新下单（用户体验最差）。
  \end{itemize}
  \item 任何业务方案都要求技术层把"看到旧价"的窗口压到\textbf{秒级以内}。
  \item 延迟成本不在 Redis，是\textbf{客服补救}（下架下单）+ \textbf{财务退差}（改价旧单）+ \textbf{用户流失}（看到旧值）。
\end{itemize}
\keypoint{延迟双删 + 短 TTL = 把"看到旧值"的窗口压到分钟以内。}
\end{frame}

% ============== 5. TTL 10min ==============
\section{TTL 10min 的业务依据}

\begin{frame}{为什么 TTL 定 10 分钟}
\begin{columns}[T,onlytextwidth]
\column{0.50\linewidth}
\begin{itemize}\small
  \item \texttt{ProductDetailTTL = 10 * time.Minute}
  \item TTL 太短：命中率低，DB 压力大。
  \item TTL 太长：商家改商品几小时还看不到。
  \item 10 min 对应：
  \begin{itemize}
    \item 商家改商品平均间隔 $\gg$ 10 min。
    \item 用户能容忍"几分钟内生效"。
    \item 单条记录 10 min 内多次访问。
  \end{itemize}
  \item TTL 只是\textbf{兜底}，主动 DEL 才是日常路径。
\end{itemize}

\column{0.50\linewidth}
\begin{tikzpicture}[scale=0.65]
  \draw[->,thick] (0,0) -- (8,0) node[right]{\small TTL};
  \draw[->,thick] (0,0) -- (0,4.5) node[above]{\small 收益};
  \draw[thick,cMain] plot[smooth] coordinates {(0.3,0.4) (1,1.4) (2,2.3) (3,2.9) (4,3.3) (5,3.5) (6,3.5) (7,3.3)};
  \draw[thick,dashed,cBad] plot[smooth] coordinates {(0.3,4.0) (1,3.7) (2,3.2) (3,2.6) (4,2.0) (5,1.4) (6,0.9) (7,0.4)};
  \node[cMain,font=\scriptsize] at (5.7,3.9) {命中率};
  \node[cBad,font=\scriptsize] at (5.9,1.7) {新鲜度};
  \draw[dotted,thick,cAccent] (3.2,0) -- (3.2,4.3);
  \node[cAccent,font=\scriptsize] at (3.2,4.5) {10 min};
\end{tikzpicture}
\end{columns}
\srcnote{repository/cache/product.go:14}
\end{frame}

% ============== 6. 延迟双删 ==============
\section{延迟双删的业务意义}

\begin{frame}[fragile]{ProductUpdate 的双删写路径}
\begin{lstlisting}[language=Go,firstnumber=271,
  caption={\srcnote{service/product.go:271--292}}]
func (s *ProductSrv) ProductUpdate(ctx context.Context, req *types.ProductUpdateReq) (resp interface{}, err error) {
    product := &model.Product{
        Name: req.Name, CategoryID: req.CategoryID,
        Title: req.Title, Info: req.Info,
        Price: req.Price, DiscountPrice: req.DiscountPrice,
        OnSale: req.OnSale,
    }
    _ = cache.DelProductDetail(ctx, req.ID)                       // 1. 第一次删
    err = dao.NewProductDao(ctx).UpdateProduct(req.ID, product)   // 2. 写 DB
    if err != nil { log.LogrusObj.Error(err); return }
    cache.DoubleDeleteAsync(req.ID, 0)                            // 3+4. 异步 500ms 后再删
    emitProductChanged(ctx, req.ID, "update")
    return
}
\end{lstlisting}
\end{frame}

\begin{frame}{为什么"再删一次"重要}
\begin{tikzpicture}[node distance=8mm]
  \foreach \name/\x in {写者 W/0, DB/30, Redis/60, 读者 R/95} {
    \node[node-box,minimum width=14mm] at (\x mm,0) (\name) {\name};
    \draw[lifeline] (\x mm,-2mm) -- (\x mm,-52mm);
  }
  \draw[arrow] (0,-8mm) -- node[above,font=\tiny]{1. DEL} (60mm,-8mm);
  \draw[arrow-bad] (95mm,-15mm) -- node[above,font=\tiny\color{cBad}]{GET miss} (60mm,-15mm);
  \draw[arrow-bad] (95mm,-22mm) -- node[above,font=\tiny\color{cBad}]{读 DB (旧值)} (30mm,-22mm);
  \draw[arrow] (0,-29mm) -- node[above,font=\tiny]{2. UPDATE} (30mm,-29mm);
  \draw[arrow-bad] (95mm,-36mm) -- node[above,font=\tiny\color{cBad}]{SET 旧值回缓存} (60mm,-36mm);
  \draw[arrow,draw=cOK] (0,-45mm) -- node[above,font=\tiny\color{cOK}]{3+4. DEL again (500ms)} (60mm,-45mm);
\end{tikzpicture}
\vspace{1mm}
\small 读者 R 在写者 W 之间穿插，把"旧值"塞回缓存。\textbf{第二次 DEL 清掉它}。
\textbf{业务感知}：旧值最多存活 500\,ms —— 客服几乎不会收到投诉。
\srcnote{repository/cache/product.go:16 \texttt{ProductDelayInterval = 500ms}}
\end{frame}

% ============== 7. 单飞回源 ==============
\section{单飞回源 (SETNX)}

\begin{frame}{SETNX 单飞：一个请求回源，其他人等结果}
\begin{tikzpicture}[node distance=5mm]
  \node[node-box] (r1) {Req 1};
  \node[node-box,below=of r1] (r2) {Req 2};
  \node[node-box,below=of r2] (r3) {Req 3};
  \node[node-box,below=of r3] (rn) {... 1\,000};
  \node[node-box,right=18mm of r2,node-warn] (lock) {SETNX\\\scriptsize product:lock};
  \node[node-box,right=18mm of lock,node-ok] (db) {DB};
  \draw[arrow] (r1.east) -- node[above,font=\tiny]{抢到} (lock);
  \draw[arrow-async] (r2.east) -- (lock);
  \draw[arrow-async] (r3.east) -- (lock);
  \draw[arrow-async] (rn.east) -- (lock);
  \draw[arrow] (lock) -- node[above,font=\tiny]{1 个回源} (db);
\end{tikzpicture}
\vspace{2mm}
\begin{itemize}
  \item \textbf{业务风险}：大主播带货 $\to$ 单个 \texttt{product\_id} 瞬间几万 QPS。缓存正好失效 $\to$ 全部打 DB $\to$ 击穿。
  \item 1 个赢家回源 DB，其他 sleep 50\,ms 重读缓存（多半命中）。
  \item \texttt{ProductLockTTL = 3s}：DB 查询 P99 < 500\,ms，3\,s 给 6 倍冗余；崩溃最多锁死 3\,s。
\end{itemize}
\srcnote{repository/cache/product.go:55--62 + service/product.go:50--58}
\end{frame}

\begin{frame}[fragile]{ProductShow 的回源锁逻辑}
\begin{lstlisting}[language=Go,firstnumber=42,
  caption={\srcnote{service/product.go:42--68}}]
cached := &types.ProductResp{}
if cacheErr := cache.GetProductDetail(ctx, req.ID, cached); cacheErr == nil {
    return cached, nil
}
locked, _ := cache.TryProductLock(ctx, req.ID)
if !locked {
    time.Sleep(50 * time.Millisecond)
    if cacheErr := cache.GetProductDetail(ctx, req.ID, cached); cacheErr == nil {
        return cached, nil
    }
} else {
    defer cache.UnlockProduct(ctx, req.ID)
}
pResp, err := s.loadProductFromDB(ctx, req.ID)
if err != nil { return nil, err }
cache.SetProductDetail(ctx, req.ID, pResp)
return pResp, nil
\end{lstlisting}
\end{frame}

% ============== 8. 压测 ==============
\section{压测数据}

\begin{frame}{商品详情 PK 查询的极限}
\begin{tabular}{@{}lrr@{}}
\toprule
链路 & RPS & p95 \\
\midrule
Ping 基线 (gin + 中间件) & 64{,}254 & 3.51\,ms \\
商品详情 PK 查询 (DB，未接缓存) & \textbf{62{,}226} & \textbf{3.01\,ms} \\
\bottomrule
\end{tabular}

\vspace{3mm}
\begin{itemize}
  \item \texttt{product/show} 走完 \texttt{ClientIP} $\to$ MySQL PK 查询的链路，跟 \texttt{/ping} \textbf{几乎重合}。
  \item 单点查询基本不增加延迟 —— 瓶颈在\textbf{热点商品并发}和\textbf{大表 list 查询}上。
  \item \textbf{业务解读}：缓存\textbf{首要价值不是延迟}，而是\textbf{削峰}（防击穿）和\textbf{容错}（DB 故障兜底）。
\end{itemize}
\keypoint{把缓存看成"业务保险"，比看成"加速器"更准确。}
\srcnote{stressTest/REPORT.md 第 36--37 行 + stressTest/product\_show.js}
\end{frame}

% ============== Q&A ==============
\appendix

\begin{frame}{面试 Q\&A（上）}
\textbf{Q1}：详情页为什么是电商最敏感的页？\\
\hspace{4mm}\small A：流量入口的终点 + 转化漏斗的最后一步；任何延迟 / 不一致直接砸 GMV。

\vspace{2mm}
\textbf{Q2}：\texttt{BossName} / \texttt{BossAvatar} 冗余是不是数据库设计错误？\\
\hspace{4mm}\small A：不是。详情页是读多写少的极致代表，冗余换的是\textbf{少一次 JOIN} + 整页缓存友好。

\vspace{2mm}
\textbf{Q3}：\texttt{Num} 和 Redis 库存的关系？\\
\hspace{4mm}\small A：\texttt{Num} 是 DB 真值（对账依据），Redis 是\textbf{下单预占运行时余量}（见 Deck 02）。
\end{frame}

\begin{frame}{面试 Q\&A（下）}
\textbf{Q4}：为什么 TTL 选 10 分钟？\\
\hspace{4mm}\small A：折中命中率和数据新鲜度；商家改商品平均间隔远 $>$ 10 min，业务能接受这个延迟。

\vspace{2mm}
\textbf{Q5}：缓存延迟带来的业务成本主要是哪几块？\\
\hspace{4mm}\small A：客服补救（下架后下单）、财务退差（改价旧单）、用户流失（看到旧值）。

\vspace{2mm}
\textbf{反追问}：商品详情单点查询不是瓶颈，为什么还要做缓存？\\
\hspace{4mm}\small A：单点不是瓶颈，但热点商品并发会击穿 DB；缓存价值是\textbf{削峰 + 容错}。
\end{frame}

\begin{frame}{代码位置一览}
\begin{description}
  \item[字段定义]\srcnote{repository/db/model/product.go:12--26}
  \item[Cache Aside 读]\srcnote{service/product.go:42--68 \texttt{ProductShow}}
  \item[延迟双删写]\srcnote{service/product.go:271--292 \texttt{ProductUpdate}}
  \item[TTL / 锁常量]\srcnote{repository/cache/product.go:13--17}
  \item[回源锁 SETNX]\srcnote{repository/cache/product.go:55--62}
  \item[DoubleDeleteAsync]\srcnote{repository/cache/product.go:64--74}
  \item[压测脚本]\srcnote{stressTest/REPORT.md 第 36--37 行 + stressTest/product\_show.js}
\end{description}
\vspace{2mm}
\keypoint{复现：\texttt{k6 run stressTest/product\_show.js}}
\end{frame}

\end{document}
